\section{Reflection-positive laws on residue tuples}\label{sec:tuples}

We now construct a probability law on finitely many residue labels at each
large prime.  The law must serve two purposes: reflection positivity will
give a multilinear H\"older inequality, while a designated cycle of labels
will have admissible increments except on a small part of the measure.
We adapt the reflection-positive construction of
\cite[Section~3]{IntersectivePower}, imposing the allowed regular argument condition
from the preceding section.  The proof below includes the positivity
argument, since the small exceptional measure and the uniformity in the
auxiliary index $l$ are both needed later.

\subsection{Slots, reflections, and the required properties}

Keep $L$ from Lemma~\ref{lem:admissible-cycle} and the exponents
$\delta$, $\gamma=\delta/10$ from \eqref{eq:local-exponents} fixed.  Choose an integer
$t\ge1$ such that $(t+1)\delta/2>5$, and set
\[
 w=2t+1,\qquad K=L!,\qquad r_*=K/2.
\]
Let $\Pi$ be the set of permutation words of length $L$: a word is a
bijection from the positions $[L]$ to the symbols $[L]$.  The slots of our
tuple are
\[
 V=\Pi^w,\qquad d=|V|=K^w.
\]
The entries of a slot are called its blocks.  Fix an $L$-cycle $c$ on the
positions.  Composition $\pi c$ therefore permutes the positions of a
word $\pi$.  Identifying $c^j$ itself with a word, fix the distinct slots
\[
 v_j=(c^j,\ldots,c^j)\quad(0\le j<L),\qquad v_L=v_0.
\]
These slots, in this order, form the designated cycle.

A \emph{cut} is specified by a block $b\in[w]$ and an ordered pair of
distinct symbols $\alpha,\beta$.  Its plus side $V_+$ consists of slots
in whose $b$th word $\alpha$ precedes $\beta$; its minus side is
$V_-=V\setminus V_+$.  The involution $\theta$ exchanges these two
symbols in block $b$ and fixes all other blocks.  Thus
$|V_+|=|V_-|=d/2$, and $\theta$ identifies the two halves.
For a measure $\nu$ on $\F_p^V$, the matrix of masses across this
cut has rows and columns indexed by $a,a'\in\F_p^{V_+}$ and entry
\[
 \nu\{z_v=a_v, z_{\theta v}=a'_v\text{ for every }v\in V_+\}.
\]
We call $\nu$ \emph{reflection positive} if every such matrix is symmetric
positive semidefinite.  Equivalently, each matrix is symmetric and
\[
 \int J((z_v)_{v\in V_+})J((z_{\theta v})_{v\in V_+})\,d\nu(z)\ge0
\]
for every real function $J$ on $\F_p^{V_+}$.

\begin{proposition}[Tuple laws]\label{prop:tuple-laws}
There is a threshold $P>P_1$, depending only on $h$, with the following
property.  For every $l\ge1$ and prime $p\ge P$, there are a probability
measure $\mu_{l,p}$ on $\F_p^V$ and a submeasure
$0\le\varepsilon_{l,p}\le\mu_{l,p}$ such that:
\begin{enumerate}
\item $\mu_{l,p}$ is reflection positive.  Both $\mu_{l,p}$ and
$\varepsilon_{l,p}$ are invariant under simultaneous translation of all
labels.  In particular, any nonzero one of these measures, or their
difference, has uniform single-slot marginals after normalization.
\item Writing
\[
 \eta_{l,p}=\varepsilon_{l,p}(1),\qquad
 \mu_{l,p}^{\circ}=\frac{\mu_{l,p}-\varepsilon_{l,p}}{1-\eta_{l,p}},
\]
we have $\eta_{l,p}\le p^{-4}$.  Under $\mu_{l,p}^{\circ}$ every
increment $z_{v_{j+1}}-z_{v_j}$, $0\le j<L$, is admissible at $l$.
\item For any integer $D\ge1$ with $p\nmid D$ and any
$a'\in\F_p$, let $A_{D,a'}$ act on all labels by
$z_v\mapsto\lambda(D)^{-1}(z_v-a')$.  Then
\begin{equation}\label{eq:tuple-covariance}
 (A_{D,a'})_*\mu_{l,p}=\mu_{lD,p},\qquad
 (A_{D,a'})_*\varepsilon_{l,p}=\varepsilon_{lD,p}.
\end{equation}
\item For $\rho_p=\mu_{l,p}$ or $\mu_{l,p}^{\circ}$, every
$v\in V$, and every $f:\F_p\to\C$ of uniform mean zero,
\begin{equation}\label{eq:tuple-mixing}
 \left\|\mathbb E_{\rho_p}
   \bigl[f(z_v)\mid(z_u)_{u\ne v}\bigr]\right\|_{L^2(\rho_p)}
 \le p^{-\gamma}\|f\|_{L^2(\F_p)}.
\end{equation}
\end{enumerate}
All thresholds and constants are independent of $l$.
\end{proposition}

The construction starts from independent labels weighted by admissible
increments.  Such a law has the cycle property and the required mixing,
but need not be reflection positive.  We add components in which selected
blocks are forgotten when labels are assigned.  For cuts in those blocks,
reflected slots share a list of labels.  We choose bounded interactions
between their selection indices so that the cut quadratic form controls
the mean square of its test function.  This lower bound will dominate the possible
negative contribution from the preceding component.
The last components need no cycle constraints, and form the exceptional
submeasure.

\subsection{Two estimates for weighted graphs}

Fix $l$ and $p\ge P_1$, and write $F=h_l$.  Define
\begin{equation}\label{eq:tuple-edge-weight}
 W(y)=\frac1k\#\{x\in\F_p:F(x)=y, x\text{ allowed regular}\},
 \qquad c_p=\mathbb E_{y\in\F_p}W(y).
\end{equation}
Since $F\bmod p$ is nonconstant of degree at most $k$, we have
$0\le W\le1$ and $c_p\ge1/(2k)$.  Lemma~\ref{lem:local-sums}
for allowed regular arguments gives, with negative-sign Fourier
coefficients,
\begin{equation}\label{eq:tuple-edge-fourier}
 |\widehat W(a)|\le p^{-4\delta}\qquad(a\in\F_p\setminus\{0\}).
\end{equation}
Indeed $\widehat W(a)$ is the normalized mean over allowed regular
arguments multiplied by their proportion divided by $k$.

\begin{lemma}[Weighted graph estimates]\label{lem:graph-estimates}
Let $G$ be a fixed finite directed graph with no loops and with at most one
edge on each unordered pair of vertices.  Give its vertices independent
uniform labels $Y_x\in\F_p$, and put
\[
 I_G(Y)=\prod_{x\to y\in E(G)}W(Y_y-Y_x),\qquad e_G=|E(G)|.
\]
Then
\begin{equation}\label{eq:tuple-graph-mean}
 \mathbb E I_G=c_p^{e_G}+O_G(p^{-4\delta}).
\end{equation}
If the vertices are partitioned into two classes and every edge joins
different classes, regard $I_G$ as a kernel between their uniform label
spaces.  Its associated operator satisfies
\begin{equation}\label{eq:tuple-graph-operator}
 \|I_G-c_p^{e_G}\|_{L^2\to L^2}=O_G(p^{-\delta}).
\end{equation}
Both bounds are uniform in $l$.
\end{lemma}

\begin{proof}
The convolution kernel $W(y-x)-c_p$ has $L^2$ operator norm at most
$p^{-4\delta}$ by \eqref{eq:tuple-edge-fourier}.  Telescope the product
$I_G$ by replacing one edge weight at a time by $c_p$.  In a resulting
term, hold fixed all labels except those at the endpoints $x,y$ of the
edge being replaced.  As there is no second edge on this pair, the other
factors depending on these two labels separate as $a(Y_x)b(Y_y)$,
with $|a|,|b|\le1$.  The convolution bound controls the conditional
average by $p^{-4\delta}$.  Summing over the edges proves
\eqref{eq:tuple-graph-mean}.

For the second assertion write $A,B$ for the label spaces of these two classes and
$D(a,b)=I_G(a,b)-c_p^{e_G}$.  In uniform $L^2$ spaces, the fourth
Schatten moment of its operator $T_D$ is
\[
 \operatorname{tr}\bigl((T_DT_D^*)^2\bigr)
 =\mathbb E_{\substack{a,a'\in A\\b,b'\in B}}
   D(a,b)D(a',b)D(a',b')D(a,b').
\]
All kernels here are real.  Expand the four differences.  Each term is
the graph average of the selected copies of $I_G$, times the corresponding
power of $-c_p^{e_G}$.  Its formal vertices consist of two separate
copies of each of the two vertex classes.  No pair of formal vertices receives two
edges: it determines one corner of the rectangle and one edge of $G$.
Thus \eqref{eq:tuple-graph-mean} applies to every term.  Replacing all its
edge weights by $c_p$ makes the expanded sum zero, so the fourth moment
is $O_G(p^{-4\delta})$.  The fourth power of the largest singular value
is at most this moment, proving \eqref{eq:tuple-graph-operator}.
\end{proof}

\subsection{Components obtained by forgetting blocks}

For $S\subseteq[w]$ with $s=|S|\le t+1$, call the blocks in $S$
marked.  For $v\in V$, its \emph{address} $u_S(v)$ is the tuple of its
words in the unmarked blocks; the address set is
$U_S=\Pi^{[w]\setminus S}$.  At every address $u$ take a list
\[
 (Y_{u,1},\ldots,Y_{u,\ell_s}),\qquad \ell_s=r_*^s,
\]
with all entries independent and uniform in $\F_p$, independently
over the addresses as well.  A slot of address $u$ will choose one of
the labels in this list.

For $s\le t$ put a directed edge $u\to u'$ between addresses whenever
\begin{equation}\label{eq:tuple-address-edge}
 u'_b=u_bc\quad\text{in at least }t+1\text{ unmarked blocks }b.
\end{equation}
Every entry of the first list is joined to every entry of the second,
with weight
\[
 I_S(Y)=\prod_{u\to u'}\prod_{i,i'\in[\ell_s]}
         W(Y_{u',i'}-Y_{u,i}).
\]
There are no loops, since $c\ne1$.  There are also no opposite edges:
because $c^2\ne1$, the witnessing blocks for $u\to u'$ and for
$u'\to u$ would be disjoint, requiring $2t+2>w$ blocks.  Consequently
the graph on individual list entries satisfies Lemma~\ref{lem:graph-estimates}.
For $s=t+1$ define $I_S=1$; these are precisely the components in which
we will dispense with the cycle constraints.

Independently select indices $j_v\in[\ell_s]$ uniformly at all slots.
To obtain the mean-square lower bound when the cut block is marked,
we put a positive interaction on some pairs of these indices.  Let
$\Delta=1/(10L^2)$.  For distinct slots $x,y$, set $H^S_{xy}=\Delta^g$
if they agree except in one marked block and the words in that block
differ by exchanging symbols at positions of distance $g$.  In all other
cases set $H^S_{xy}=0$, and also set $H^S_{xx}=0$.  Define
\begin{align}
 R_S(j)&=\exp\left(\sum_{\{x,y\}\subseteq V}
                      H^S_{xy}\mathbf1_{\{j_x=j_y\}}\right),
 \label{eq:tuple-index-weight}\\
 \nu_S(\Phi)&=\mathbb E_{Y,j}
       I_S(Y)R_S(j)\,
       \Phi\bigl((Y_{u_S(v),j_v})_{v\in V}\bigr).
 \label{eq:tuple-component}
\end{align}
The sum in \eqref{eq:tuple-index-weight} is over unordered two-element
subsets.  Each $\nu_S$ is a finite nonnegative measure, with mass bounded
above by a fixed constant: $I_S\le1$ and $1\le R_S\le C$ for a constant
depending only on the fixed slot set.

Every component is invariant under each reflection $\theta$.  To see
this, exchanging two symbols throughout one block preserves all
transposition distances used in $H^S$.  For an unmarked block, the same
symbol exchange permutes addresses and preserves
\eqref{eq:tuple-address-edge}, since symbol permutations commute with
the right action of $c$ on positions.  For a marked block, addresses are
unchanged.  The uniform lists and indices are preserved in either case.
Thus the cut matrices of $\nu_S$ are symmetric.  Their positivity is
the remaining issue.

\subsection{Positive contributions and their comparison}

Fix a cut in block $b$, and for a real function $J$ of the plus labels
write its quadratic form as
\[
 Q_S(J)=\int J((z_v)_{v\in V_+})
             J((z_{\theta v})_{v\in V_+})\,d\nu_S(z).
\]
We first prove that marking $b$ gives a strictly positive lower bound.
We then compare that lower bound with the possible negative contribution
when $b$ is unmarked.  These two estimates explain both the list lengths
$\ell_s$ and the coefficients in the final mixture.

Suppose first that $b\in S$.  The matrix of interactions across the cut,
indexed by $x,y\in V_+$, is
\[
 \mathsf A_{xy}=H^S_{x,\theta y}.
\]
It is symmetric by reflection invariance.  If the two cut symbols in
slot $x$ occupy positions of distance $g$, then
$\mathsf A_{xx}=\Delta^g$.  Any other transposition taking $x$ across
the cut must move one cut symbol past the other, leaving the latter
fixed.  It therefore exchanges positions at distance at least $g+1$.
Only a transposition in block $b$ can do this.  There are fewer than
$L^2$ choices, whence
\[
 \sum_{y\ne x}|\mathsf A_{xy}|\le L^2\Delta^{g+1}
       =\tfrac1{10}\Delta^g.
\]
By diagonal dominance,
$\mathsf A\succeq\kappa\operatorname{Id}$ with the fixed choice
$\kappa=(9/10)\Delta^{L-1}>0$.

Here is why this positivity for interactions gives strict positivity for
the whole index kernel.  Identify minus indices with plus indices using
$\theta$, and let $j,j'$ be two half-index assignments.  Put
\[
 u(j)=(\mathbf1_{\{j_x=i\}})_{x\in V_+,\ i\in[\ell_s]}.
\]
The cross factor of $R_S$ is the matrix with entries
\[
 C(j,j')=\exp\bigl(u(j)^T(\mathsf A\otimes\operatorname{Id})u(j')\bigr).
\]
Split $\mathsf A$ as $\kappa\operatorname{Id}$ plus a positive
semidefinite matrix.  The entrywise exponential of a positive
semidefinite Gram matrix is positive semidefinite: expand its power
series and apply the Schur product theorem to each power.  The factor
from $\mathsf A-\kappa\operatorname{Id}$ therefore is positive
semidefinite and has diagonal entries at least one.  The factor from
$\kappa\operatorname{Id}$ is the tensor product, over $x\in V_+$, of
\[
 \mathbf1\mathbf1^T+(e^\kappa-1)\operatorname{Id}.
\]
It dominates $(e^\kappa-1)^{d/2}\operatorname{Id}$.  Taking its Schur
product with the other factor preserves this lower bound, since the
Schur product with the identity is the diagonal of that factor.
The internal interactions in the two halves multiply this matrix on
the left and right by the same diagonal matrix with entries at least one,
so the lower bound is still valid.

Reflected slots have the same address when $b\in S$.  For fixed lists
$Y$, the two factors involving $J$ consequently evaluate the same vector
indexed by half-index assignments.  There are
$n_s=\ell_s^{d/2}$ such assignments.  The uniform average over both
halves is $n_s^{-2}$ times the associated quadratic form; its diagonal
lower bound is thus $(e^\kappa-1)^{d/2}n_s^{-1}$ times the uniform
average of the squares.  Since there are only finitely many possible
$s$, we obtain a fixed $c_*>0$ such that
\begin{equation}\label{eq:tuple-marked-positive}
 Q_S(J)\ge c_*\,
 \mathbb E_{Y,j_+}I_S(Y)
    J((Y_{u_S(v),j_v})_{v\in V_+})^2\qquad(b\in S).
\end{equation}

Now suppose $b\notin S$.  The lists split into plus and minus addresses
according to the order of the cut symbols in block $b$.  Write $Y_+$
for the plus lists, and $I_+(Y_+)$ for the factors of $I_S$ whose edges
have both ends on that side.  Set
\[
 \mathcal A_S(J)=\mathbb E_{Y_+,j_+} I_+(Y_+)
       J((Y_{u_S(v),j_v})_{v\in V_+})^2.
\]
No index interaction crosses this cut, because changing a marked block
does not change the word in block $b$.  Let $R_+$ be the internal index
weight on the plus side.  After averaging the indices, $Q_S(J)$ is the
quadratic form of the kernel of cross-list edge weights applied to the
half-function
\[
 B(Y_+)=I_+(Y_+)\,
   \mathbb E_{j_+}R_+(j_+)J((Y_{u_S(v),j_v})_{v\in V_+}).
\]
By Cauchy--Schwarz, $R_+\le C$, and $I_+^2\le I_+$,
$\|B\|_2^2\le C\mathcal A_S(J)$.
Lemma~\ref{lem:graph-estimates} replaces the cross kernel by its
nonnegative constant value with operator error $O(p^{-\delta})$.
The constant kernel contributes a nonnegative multiple of
$(\mathbb E B)^2$.  Hence, uniformly over $S$ and the cut,
\begin{equation}\label{eq:tuple-unmarked-error}
 Q_S(J)\ge-Cp^{-\delta}\mathcal A_S(J)\qquad(b\notin S).
\end{equation}
If $s=t+1$, the cross kernel is exactly $1$, and $Q_S(J)\ge0$.

For $b\notin S$ and $s\le t$, let $S'=S\cup\{b\}$.  We claim that
\begin{equation}\label{eq:tuple-concatenation}
 Q_{S'}(J)\ge c_*r_*^{-d/2}\mathcal A_S(J).
\end{equation}
The new address forgets block $b$.  Exactly $r_*=K/2$ old plus
addresses become each new address, one for each word in which $\alpha$
precedes $\beta$.  Enumerate those words as $\sigma_1,\ldots,\sigma_{r_*}$,
and write the old addresses above a new address $u$ as $(u,\sigma_a)$.
Concatenate their old lists by setting
\[
 \widetilde Y_{u,(a-1)\ell_s+i}=Y_{(u,\sigma_a),i}
 \qquad(a\in[r_*],\ i\in[\ell_s]).
\]
These are independent uniform lists of length
$r_*\ell_s=\ell_{s+1}$, precisely as required by the definition of
$\nu_{S'}$.

Under this identification of list variables, $I_{S'}\ge I_+$ pointwise.
Indeed, any edge between new addresses is witnessed by at least $t+1$
blocks other than $b$, so every pair of their constituent old plus
addresses already had that directed edge.  Thus every factor of
$I_{S'}$ occurs in $I_+$; the latter may have additional factors in
$[0,1]$.  If $|S'|=t+1$, the same inequality follows from $I_{S'}=1$.
In the nonnegative right side of \eqref{eq:tuple-marked-positive} for
$S'$, denote the independently chosen uniform new indices by $j'_v$.
For a plus slot of old address $(u,\sigma_a)$, restrict $j'_v$ to
$\{(a-1)\ell_s+1,\ldots,a\ell_s\}$.  Each restriction has probability
$1/r_*$, independently over the $d/2$ slots.  Conditional on these
restrictions, the selected labels have exactly their old distribution.
Using $I_{S'}\ge I_+$ gives \eqref{eq:tuple-concatenation}.
Figure~\ref{fig:list-concatenation} illustrates this identification of
the lists and the index restrictions.

\begin{figure}[!bp]
  \centering
  \begin{tikzpicture}[
      x=1cm,y=1cm,>=Stealth,
      every node/.style={font=\small},
      oldlist/.style={draw=black!65,rounded corners=1pt,minimum height=8mm,
        minimum width=32mm,inner sep=2pt},
      guide/.style={draw=black!65,->,line width=.5pt},
      chosen/.style={fill=blue!7,draw=blue!55!black},
      index/.style={text=blue!55!black}
    ]
    \node[anchor=south] at (5.4,.9)
      {Old plus lists whose addresses forget to $u$};
    \node[anchor=south] at (1.6,.42) {$(u,\sigma_1)$};
    \node[anchor=south] at (5.4,.42) {$(u,\sigma_a)$};
    \node[anchor=south] at (9.2,.42) {$(u,\sigma_{r_*})$};
    \node[oldlist] (first) at (1.6,0)
      {$(Y_{(u,\sigma_1),i})_{i\in[\ell_s]}$};
    \node[oldlist,chosen] (middle) at (5.4,0)
      {$(Y_{(u,\sigma_a),i})_{i\in[\ell_s]}$};
    \node[oldlist] (last) at (9.2,0)
      {$(Y_{(u,\sigma_{r_*}),i})_{i\in[\ell_s]}$};
    \node at (3.5,0) {$\cdots$};
    \node at (7.3,0) {$\cdots$};
    \draw[guide] (first.south) -- (1.6,-1.18);
    \draw[guide] (middle.south) -- (5.4,-1.18);
    \draw[guide] (last.south) -- (9.2,-1.18);
    \node[fill=white,inner sep=2pt] at (5.4,-.8) {concatenate};

    \fill[blue!7] (3.8,-1.2) rectangle (7,-2);
    \draw[black!65] (0,-1.2) rectangle (10.8,-2);
    \draw[black!65] (3.2,-1.2) -- (3.2,-2);
    \draw[black!65] (3.8,-1.2) -- (3.8,-2);
    \draw[black!65] (7,-1.2) -- (7,-2);
    \draw[black!65] (7.6,-1.2) -- (7.6,-2);
    \node at (1.6,-1.6) {sublist $1$};
    \node at (3.5,-1.6) {$\cdots$};
    \node[index] at (5.4,-1.6) {sublist $a$};
    \node at (7.3,-1.6) {$\cdots$};
    \node at (9.2,-1.6) {sublist $r_*$};
    \draw[decorate,decoration={brace,mirror,amplitude=4pt}]
      (0,-2.12) -- (10.8,-2.12);
    \node[anchor=north] at (5.4,-2.33)
      {$\widetilde Y_u:\quad \ell_{s+1}=r_*\ell_s\ \text{independent uniform entries}$};

    \node[align=center,text width=11.4cm,anchor=north] at (5.4,-3.02)
      {For $v\in V_+$ with $u_S(v)=(u,\sigma_a)$, retain only indices\\[3pt]
       $\displaystyle j'_v\in\{(a-1)\ell_s+1,\ldots,a\ell_s\},
       \qquad \mathbb P(\text{retained at every }v\in V_+)=r_*^{-d/2}.$};
  \end{tikzpicture}
  \caption{Marking the cut block $b$: the comparison between $S$ and
    $S'=S\cup\{b\}$. For each new address $u$, concatenate its $r_*$ old
    plus lists; the displayed sublists are schematic. Independence and
    probabilities refer to sampling before weights are applied.
    Restricting the new index at each plus slot to its former sublist
    recovers the old distribution of selected labels, with probability
    $r_*^{-d/2}$. Every edge between new addresses
    gives an old edge between every pair of their plus sublists. Hence every
    factor of $I_{S'}$ occurs in $I_+$, and $0\le W\le1$ gives
    $I_{S'}\ge I_+$. These are the two ingredients of the comparison for
    $Q_{S'}(J)$.}
  \label{fig:list-concatenation}
\end{figure}


\subsection{The mixture and its cycle property}

Define a finite measure and its normalization by
\begin{equation}\label{eq:tuple-mixture}
 M=\sum_{\substack{S\subseteq[w]\\|S|\le t+1}}
           p^{-|S|\delta/2}\nu_S,
 \qquad \mu_{l,p}=\frac{M}{M(1)}.
\end{equation}
For a fixed cut in block $b$, pair each $S$ missing $b$, $|S|\le t$,
with $S'=S\cup\{b\}$.  By
\eqref{eq:tuple-unmarked-error} and \eqref{eq:tuple-concatenation}, their
combined quadratic form is at least
\[
 p^{-s\delta/2}
 \bigl(c_*r_*^{-d/2}p^{-\delta/2}-Cp^{-\delta}\bigr)
 \mathcal A_S(J)\ge0
\]
for all sufficiently large $p$.  Every component containing $b$ is
used in exactly one such pair; the remaining components have
$|S|=t+1$, $b\notin S$, and nonnegative quadratic forms.  This proves
reflection positivity of $M$ and hence of $\mu_{l,p}$.

For $S=\varnothing$, every list has length one, all index weights are
one, and the component mass is a graph average.  Therefore
\eqref{eq:tuple-graph-mean} and $c_p\ge1/(2k)$ show that
$\nu_\varnothing(1)\ge c_0>0$ for sufficiently large $p$, uniformly in
$l$.  The uniform upper bounds on all component masses give
\[
 c_0\le M(1)\le C_0.
\]
Define the exceptional submeasure by
\begin{equation}\label{eq:tuple-exceptional}
 \varepsilon_{l,p}
 =\frac{p^{-(t+1)\delta/2}}{M(1)}
       \sum_{|S|=t+1}\nu_S.
\end{equation}
It has mass $O(p^{-(t+1)\delta/2})\le p^{-4}$ after increasing the
prime threshold, by the choice of $t$.

For $|S|\le t$ at least $t+1$ blocks remain unmarked, and the addresses
of $v_j,v_{j+1}$ satisfy \eqref{eq:tuple-address-edge} in every one of
them.  The product $I_S$ therefore includes the weight on every pair of
their list entries.  On its support, in particular, the entries selected
at those two slots have increment in the support of $W$, which consists
of admissible residues.  This proves the cycle property for
$\mu_{l,p}^{\circ}$.  Translating every list entry by the same element
of $\F_p$ leaves all the weights unchanged.  Each component,
and hence both measures in the proposition, is translation invariant.
The assertion about one-slot marginals follows because translation acts
transitively on $\F_p$.

\subsection{Covariance and conditional mixing}

It remains to show that the construction follows the auxiliary family
under affine changes and that conditioning on the other slots reveals
little about a mean-zero function at one slot.  The first property will
permit passage to progressions.  The second will control large
denominators in the Fourier lifts.

If $p\nmid D$, the auxiliary-family identity \eqref{eq:auxiliary-change} gives
\[
 W_{lD,p}(y)=W_{l,p}(\lambda(D)y).
\]
Indeed the substitution $x\mapsto Dx-i$ is bijective modulo $p$,
preserves $T_l(Dx-i)=T_{lD}(x)$, and preserves nonvanishing of the
derivative because both $D$ and $\lambda(D)$ are units modulo $p$.
Consequently changing every list entry to
$\lambda(D)^{-1}(Y_{u,i}-a')$ transforms each component for $l$ into
the component for $lD$, preserving its mass and all index weights.
The mixtures and their exceptional parts therefore satisfy
\eqref{eq:tuple-covariance}.

For \eqref{eq:tuple-mixing}, first consider the normalized component
$\nu_\varnothing/\nu_\varnothing(1)$.  At a fixed slot $v$ write
$x=z_v$ and $y=(z_u)_{u\ne v}$.  Its unnormalized density, relative to
independent uniform labels, factors as
\[
 J_1(x,y)J_2(y),
\]
where $J_1$ is the product of the $m$ edge weights incident with $v$,
and $J_2$ contains the remaining edges.  Put
\[
 A(y)=\mathbb E_xJ_1(x,y),\qquad
 B(y)=\mathbb E_xf(x)J_1(x,y).
\]
Apply \eqref{eq:tuple-graph-operator} to the graph of incident edges,
whose two sides are $\{v\}$ and the other vertices.  Since
$\mathbb E f=0$, it gives
\[
 \|A-c_p^m\|_2\le Cp^{-\delta},\qquad
 \|B\|_2\le Cp^{-\delta}\|f\|_2.
\]
The conditional expectation at $y$ is $B(y)/A(y)$ when $A(y)>0$.
The marginal density of $y$ is
$J_2(y)A(y)/\nu_\varnothing(1)$, so the squared conditional norm is
\begin{equation}\label{eq:tuple-leading-conditional}
 \frac1{\nu_\varnothing(1)}
    \mathbb E_yJ_2(y)\frac{|B(y)|^2}{A(y)},
\end{equation}
with the ratio defined as zero when $A=0$.
On $A\ge c_p^m/2$, this is bounded using the estimate for $B$.
The complementary set has uniform probability $O(p^{-2\delta})$
by the estimate for $A$, since $c_p^m$ is bounded below.  On that set,
weighted Cauchy--Schwarz and $J_1\le1$ give
\[
 \frac{|B(y)|^2}{A(y)}
 \le\mathbb E_x|f(x)|^2J_1(x,y)\le\|f\|_2^2.
\]
Together with $J_2\le1$ and the lower bound on
$\nu_\varnothing(1)$, this proves that
\eqref{eq:tuple-leading-conditional} is
$O(p^{-2\delta})\|f\|_2^2$.

Both $\mu_{l,p}$ and $\mu_{l,p}^{\circ}$ are mixtures of the normalized
components, and in either mixture the total probability of
$S\ne\varnothing$ is $O(p^{-\delta/2})$.  To justify transfer of the
estimate, retain the component label as an additional random variable.
Conditional Jensen shows that conditioning on this label as well as on
all other slots can only increase the squared conditional norm.  In each
normalized component the trivial bound is $\|f\|_2^2$, since that
component has a uniform marginal at slot $v$.  Averaging the leading
estimate and these trivial bounds therefore yields
\[
 \left\|\mathbb E_{\rho_p}
   [f(z_v)\mid(z_u)_{u\ne v}]\right\|_2^2
 \le C p^{-\delta/2}\|f\|_2^2.
\]
As $2\gamma=\delta/5<\delta/2$, increasing the prime threshold makes
this at most $p^{-2\gamma}\|f\|_2^2$.  This proves
\eqref{eq:tuple-mixing} and completes the proof of
Proposition~\ref{prop:tuple-laws}.

Fix henceforth a threshold $P>P_1$ for which the proposition holds, large
enough also that
\begin{equation}\label{eq:prime-thresholds}
 p^{\gamma/2}\ge3\quad(p\ge P),\qquad
 \prod_{p\ge P}(1+p^{-3})-1\le\tfrac14,
\end{equation}
where the product is over primes.  Its tail tends to one because
$\sum_p p^{-3}<\infty$.  Define
\begin{equation}\label{eq:small-prime-modulus}
 M_0=\prod_{p<P}p^{E_p}\ge2.
\end{equation}
The previously chosen $L$ is unchanged when $P$ is enlarged.  For
$D=M_0$ times a product of distinct primes at least $P$, a residue
modulo $D$ is called admissible at $l$ if it is admissible at every local
factor $p^{E_p}$ of $D$.
